
%\needspace{4\baselineskip}
\color{uniblue}{\subsection{Контроль силового выключателя (КСВ)}}
\color{black}

\begin{enumerate}[label=\arabic{section}.\arabic{subsection}.\arabic{enumi}, labelsep=4pt, leftmargin=0pt, itemindent=57pt, itemsep=0pt, parsep=0pt] % если что можно крутить itemsep

\item
Функциональный блок контроля силового выключателя (КСВ) предназначен для обработки технологической информации от выключателя и его цепей управления, а также от трансформатора тока, позволяя определять их состояние, информировать о возможных неисправностях и формировать соответствующую сигнализацию.

\item
В случае присоединения линии к системе через два выключателя для каждого из них можно задействовать свою логику КСВ. Далее приведено описание КСВ для выключателя В1, логика КСВ для второго выключателя аналогична.

\item
Контроль готовности привода выключателя осуществляется по внешнему сигналу \textit{«Пруж. не заведены В1»}. Сигнал о незаведенном состоянии пружины \textit{«Ср. Пруж. не заведены В1»} формируется с выдержкой времени \textbf{«Тср Пруж. не заведены В1»}. 

\item
Контроль целостности цепей управления осуществляется по сигналам состояния цепей первой и второй группы электромагнитов отключения (ЭМО1 и ЭМО2) и цепи электромагнита включения (ЭМВ). Наличие второй группы электромагнитов отключения определяется положением программного переключателя \textbf{«Кол-во ЭМО В1»}. При одновременном наличии или отсутствии сигналов контроля цепей ЭМО и ЭМВ в течение регулируемой выдержки времени \textbf{«Тср Неиспр. ЦУ В1»} формируется сигнал \textit{«Неиспр. ЦУ В1»}. Выдержка времени должна быть не менее времени взвода пружины привода выключателя.

\item
Функциональная схема логики контроля В1 приведена на рисунке \ref{pict_ksv:KSV_control_TT_B1}. Функциональная схема логики контроля В2 выполняется аналогично схеме логики контроля В1.

\vspace{2mm}
\begin{figure}[H]
\centering
\includegraphics[width=1\textwidth,height=1\textheight,keepaspectratio]{\fbpath/761.3371 КСВ/img/KSV_control_TT_B1.pdf}
\captionof{figure}{Функциональная схема контроля В1}\label{pict_ksv:KSV_control_TT_B1}
\end{figure}

\item
Предусмотрен контроль давления элегаза в баке выключателя по внешним сигналам \textit{«Низк. давл. SF6 В1»} и \textit{«Авар. давл. SF6 В1»} от специальных датчиков давления. Соответствующие сигналы о низком и аварийном давлении элегаза \textit{«Ср. Низк. давл. SF6 В1»} и \textit{«Ср. Авар. давл. SF6 В1»} формируются с выдержками времени, равными 1 с.

\item
Контроль положения автомата шинок питания осуществляется по внешнему сигналу \textit{«Автомат ШП В1»}. Сигнал об отключенном положении автомата \textit{«Ср. Автомат ШП В1»} формируется с выдержкой времени выдержкой времени, равной 1 с. 

\item
Обогрев выключателя контролируется по внешнему сигналу \textit{«Неиспр. обогр. В1»}. Соответствующий сигнал о неисправности обогрева \textit{«Ср. Неиспр. обогр. В1»} формируется с регулируемой выдержкой времени \textbf{«Тср Неиспр. обогр. В1»}.

\item
Сигнализация о неисправности выключателя (сигнал \textit{«Неиспр. В1»}) срабатывает при наличии одного из сигналов: \textit{«Неиспр. ЦУ В1»}, \textit{«Ср. Пруж. не заведены В1»}, \textit{«Ср. Низк. давл. SF6 В1»}, \textit{«Ср. Авар. давл. SF6 В1»}, \textit{«Ср. Автомат ШП В1»}, \textit{«Ср. Неиспр. обогр. В1»}, а также от сигнала \textit{«Ср. ЗНФ В1»}.

\item
В устройстве также предусмотрен контроль состояния ТТ. Контроль давления элегаза в ТТ осуществляется по внешним сигналам \textit{«Низк. давл. SF6 ТТ В1»} и \textit{«Авар. давл. SF6 ТТ В1»}. Соответствующие сигналы о низком и аварийном давлении элегаза \textit{«Ср. Низк. давл. SF6 В1»} и \textit{«Ср. Авар. давл. SF6 В1»} формируются с выдержками времени, равными 1 с. Обогрев ТТ контролируется по внешнему сигналу \textit{«Неиспр. обогр. ТТ В1»}. Соответствующий сигнал о неисправности обогрева \textit{«Ср. Неиспр. обогр. ТТ В1»}
формируется с регулируемой выдержкой времени \textbf{«Тср Неиспр. обогр. ТТ В1»}. По наличию одного из указанных сигналов формируется сигнал о неисправности ТТ \textit{«Неиспр. ТТ В1»}. При вводе программного ключа \textbf{«Откл. от Авар. давл. SF6 ТТ В1»} предусматривается действие на отключение выключателя при аварийном снижении давления элегаза в ТТ.

\item
Функциональная схема логики контроля ТТ В1 приведена на рисунке \ref{pict_ksv:KSV_control}. Функциональная схема логики контроля ТТ В2 выполняется аналогично схеме логики контроля ТТ В1.

\vspace{2mm}
\begin{figure}[H]
\centering
\includegraphics[width=1\textwidth,height=1\textheight,keepaspectratio]{\fbpath/761.3371 КСВ/img/KSV_control.pdf}
\captionof{figure}{Функциональная схема контроля ТТ В1}\label{pict_ksv:KSV_control}
\end{figure}


\item
Контроль наличия питания цепей сигнализации выключателя и ТТ осуществляется по внешним сигналам \textit{«Контр. опертока ЭМВ/ЭМО1 В1»}, \textit{«Контр. опертока ЭМО2 В1»} и \textit{«Контр. опертока сигн. В1»}. При состоянии «0» одного из сигналов срабатывает сигнализация о неисправности цепей сигнализации выключателя и ТТ (сигнал \textit{«Неиспр. сигн. В1/ТТ В1»}) с выдержкой времени 1 с. 

\item
Функциональная схема логики контроля опертока В1 приведена на рисунке \ref{pict_ksv:KSV_control_opercurrent}. Функциональная схема логики контроля опертока В2 выполняется аналогично схеме логики контроля опертока В1.

\vspace{2mm}
\begin{figure}[H]
\centering
\includegraphics[width=1\textwidth,height=1\textheight,keepaspectratio]{\fbpath/761.3371 КСВ/img/KSV_control_opercurrent.pdf}
\captionof{figure}{Функциональная схема контроля опертока В1}\label{pict_ksv:KSV_control_opercurrent}
\end{figure}

\item
Для электромагнитов управления предусмотрена защита от длительного протекания тока. Контроль протекания тока через ЭМВ, ЭМО1 и ЭМО2 осуществляется при помощи токовых реле, положения контактов которых должны быть заведены на соответствующие входные сигналы ФСУ «Работа ЭМВ В1», «Работа ЭМО1 В1», «Работа ЭМО2 В1». При длительном протекании тока по цепям ЭМО1 и ЭМВ формируется сигнал «Обест. ЭМО1 и ЭМВ В1», действующий на независимый расцепитель автомата питания цепей ЭМО1 и ЭМВ.

\item
Аналогично при длительном протекании тока по цепи ЭМО2 формируется сигнал \textit{«Обест. ЭМО2 В1»} для действия на независимый расцепитель автомата питания цепи ЭМО2. Регулируемые выдержки времени \textbf{«Тср защиты ЭМВ В1»}, \textbf{«Тср защиты ЭМО1 В1»} и \textbf{«Тср защиты ЭМО2 В1»} определяют время срабатывания на обесточивание ЭМВ, ЭМО1 и ЭМО2 соответственно. Кроме того, наличие сигнала \textit{«Ср. Авар. давл. SF6 В1»} действует на обесточивание всех электромагнитов.

\item
Функциональная схема логики защиты ЭМУ В1 приведена на рисунке \ref{pict_ksv:KSV_EMU_protection}. Функциональная схема логики защиты ЭМУ В2 выполняется аналогично схеме логики защиты ЭМУ В1.

\vspace{2mm}
\begin{figure}[H]
\centering
\includegraphics[width=1\textwidth,height=1\textheight,keepaspectratio]{\fbpath/761.3371 КСВ/img/KSV_EMU_protection.pdf}
\captionof{figure}{Функциональная схема защиты ЭМУ В1}\label{pict_ksv:KSV_EMU_protection}
\end{figure}

\item 
В устройстве предусмотрена логика формирования блокировки команд на включение и отключение выключателя.

\item
На блокировку отключения выключателя действует сигнал \textit{«Ср. Авар. давл. SF6 В1»}. На блокировку включения действуют сигналы \textit{«Ср. Авар. давл. SF6 В1»}, \textit{«Неиспр. ЦУ В1»}, \textit{«Ср. Пруж. не заведены В1»}, \textit{«Ср. Автомат ШП В1»}.

\item
Также предусмотрена блокировка отключения от внешнего сигнала \textit{«Внеш. блок. откл. В1»}, блокировка включения от внешнего сигнала \textit{«Внеш. блок. вкл. В1»} и блокировка управления (как включения, так и отключения) от внешнего сигнала \textit{«Внеш. блок. упр. В1»}.

\item
Функциональная схема логики блокировки управления В1 приведена на рисунке \ref{pict_ksv:KSV_control_block}. Функциональная схема логики блокировки управления В2 выполняется аналогично схеме логики блокировки управления В1.

\vspace{2mm}
\begin{figure}[H]
\centering
\includegraphics[width=1\textwidth,height=1\textheight,keepaspectratio]{\fbpath/761.3371 КСВ/img/KSV_control_block.pdf}
\captionof{figure}{Функциональная схема блокировки управления В1}\label{pict_ksv:KSV_control_block}
\end{figure}

\item
В устройстве предусмотрена программная фиксация команд или включенного положения выключателя в зависимости от положения программного переключателя \textbf{«Реле фиксации В1»}. В положении \textit{«РФК»} фиксируется команда включения, в положении \textit{«РФП»} – включенное положение. Квитирование сигнала фиксации \textit{«РФ В1»} может осуществляться от сброса сигнализации или подачей оперативной команды на отключение. При отключении выключателя от защит и наличии сигнала \textit{«РФ В1»} формируется сигнал о факте аварийного отключения выключателя \textit{«Сигнал несоотв. В1»}, необходимый для работы АПВ.

\item
Функциональная схема логики реле фиксации В1 приведена на рисунке \ref{pict_ksv:KSV_relay_fix_B1}. Функциональная схема логики реле фиксации В2 выполняется аналогично схеме логики реле фиксации В1.

\vspace{2mm}
\begin{figure}[H]
\centering
\includegraphics[width=1\textwidth,height=1\textheight,keepaspectratio]{\fbpath/761.3371 КСВ/img/KSV_relay_fix_B1.pdf}
\captionof{figure}{Функциональная схема логики реле фиксации В1}\label{pict_ksv:KSV_relay_fix_B1}
\end{figure}

\end{enumerate}